% GENERATED FILE. Edit canonical lessons, then run npm run generate:books.
% canonical-source-hash: fnv1a64:7488bd8c935f62d7
% canonical-lessons: ML-C10-azhcha

\chapter{Days of the Week}
\label{ch:ml-days-of-the-week}

\begin{chapteropening}
\textbf{By the end of this chapter:} \emph{I can name the seven days of the Malayalam week and say which planet each one is named for.}

Name the seven planet-days, separate the native ones from the Sanskrit ones, and explain their near-identity with Tamil's.
\end{chapteropening}

\section[\ml{തിങ്കൾ} \ml{ചൊവ്വ} \ml{ബുധൻ} \ml{വ്യാഴം} \ml{വെള്ളി} \ml{ശനി} \ml{ഞായർ}]{\ml{തിങ്കൾ} to \ml{ഞായർ} — the week Malayalam shares with Tamil}

\label{lesson:ML-C10-azhcha}



\begin{quote}
Malayalam and Tamil are close cousins — they were once, long ago, the same language — and nowhere does that show more clearly than in the words for the days of the week.
\end{quote}

\subsection*{What to know first}
\noindent
\begin{tabularx}{\linewidth}{@{}>{\raggedright\arraybackslash}X>{\raggedright\arraybackslash}X>{\raggedright\arraybackslash}X>{\raggedright\arraybackslash}X@{}}
\toprule
\textbf{Malayalam} & \textbf{planet-word} & \textbf{meaning} & \textbf{compare Tamil} \\
\midrule
\textbf{\ml{തിങ്കളാഴ്ച}} & \emph{thiṅkaḷ} & \textquotedblleft{}\textbf{Moon}\textquotedblright{} & \emph{thiṅgaḷ} — nearly identical \\
\textbf{\ml{ചൊവ്വാഴ്ച}} & \emph{covva} & traditionally \textquotedblleft{}\textbf{the red one}\textquotedblright{} & \emph{cevvāy} — the same word, worn differently \\
\textbf{\ml{ബുധനാഴ്ച}} & \emph{budhan} & \textbf{Mercury} & Sanskrit \textbf{Budha}, same as Tamil \emph{pudhaṉ} \\
\textbf{\ml{വ്യാഴാഴ്ച}} & \emph{vyāzham} & \textbf{Jupiter} & matches Tamil's \emph{viyāzhaṉ} closely \\
\textbf{\ml{വെള്ളിയാഴ്ച}} & \emph{veḷḷi} & traditionally \textquotedblleft{}\textbf{silver}\textquotedblright{} & \emph{veḷḷi} — the \textbf{same word} \\
\textbf{\ml{ശനിയാഴ്ച}} & \emph{śani} & \textbf{Saturn} & Sanskrit \textbf{Śani}, same as Tamil \emph{saṉi} \\
\textbf{\ml{ഞായറാഴ്ച}} & \emph{ñāyar} & \textquotedblleft{}\textbf{Sun}\textquotedblright{} & \emph{ñāyiṟu} — the same root \\
\bottomrule
\end{tabularx}

\begin{culture}
This is the closest family match you'll see across this whole course: \textbf{thiṅkaḷ/thiṅgaḷ}, \textbf{veḷḷi/veḷḷi}, and \textbf{ñāyar/ñāyiṟu} are essentially \textbf{the same words}, because Malayalam only became a distinct language from Tamil within roughly the last thousand years or so. Where Tamil says \textbf{\ta{கிழமை}} (\emph{kizhamai}) for \textquotedblleft{}day,\textquotedblright{} Malayalam says \textbf{\ta{ஆழ்ச}} (\emph{āzhcha}) — a different word for the same idea (the two aren't confirmed cognates; treat this as two languages solving the same naming problem separately, not necessarily one shared root). Like Tamil, Malayalam mixes \textbf{native} planet-words (Moon, Mars, Venus, Sun) with \textbf{Sanskrit} ones (Mercury, Jupiter, Saturn) — the same split, inherited from the same shared ancestor language.
\end{culture}

\subsection*{Your turn}
Say these aloud:

\begin{itemize}
\raggedright
  \item \textquotedblleft{}thiṅkaḷ\textquotedblright{} — moon — then \textquotedblleft{}āzhcha,\textquotedblright{} day/week
  \item the shared words — \textquotedblleft{}veḷḷi\textquotedblright{} (Venus), same in both Malayalam and Tamil
  \item the Sanskrit ones — \textquotedblleft{}budhan, vyāzham, śani\textquotedblright{} — Mercury, Jupiter, Saturn
\end{itemize}

\begin{tcolorbox}[breakable,colback=teal!4,colframe=teal!35!black,title={Before you move on}]
Why do Malayalam and Tamil's day-names look so alike? (\textbf{They were the same language until roughly a thousand years ago.}) Which planet-names does Malayalam share almost word-for-word with Tamil? (\textbf{thiṅkaḷ/Moon, veḷḷi/Venus, ñāyar/Sun} — and \emph{covva}/Mars, a worn variant of \emph{cevvāy}.) Which three come from Sanskrit in both languages? (\textbf{Mercury (Budha), Jupiter (vyāzham/viyāzhaṉ), Saturn (Śani).})
\end{tcolorbox}
